% Dies ist die Titelseite des Feinentwurfs.
% Die in "<...>" sind zu ersetzen
% Die Ausgabe darf 1 Seite nicht überschreiten, also ggf. Abstände anpassen
% Die Angabe in [...] gibt den Abstand nach der entsprechenden Zeile an.

%----Stil dieser Seite----------------------------------------------------------
\thispagestyle{plain}      % Kopfzeile bleibt leer

\begin{titlepage}

%----zentrierte Ausrichtung über die gesamte Seite----------------------------
\begin{center}

%----Titel des Praktikum (\praktikumTitel in newComments zu verändern)--------
{\relsize{4}{\textbf{\textsc{\praktikumTitel}}}}\\[5ex]

%----Titel des Teilprojektes (\projektTitel in newComments verändern)---------
{\relsize{3}{\textbf{\textsc{\projektTitel}}}}\\[5ex]

Software-Entwicklungspraktikum (SEP)\\
Sommersemester 2012\\[6ex]

{\relsize{3}\so{\textbf{Feinentwurf}}}\\[5ex]

%----eingebundenes Logo der TU--------------------------------------------------
\includegraphics[scale=0.8]{bilder/carolo.jpg}\\[5ex]

%----Daten des Auftraggebers
Auftraggeber\\
Technische Universität Braunschweig\\
<Name des Instituts>\\
<Name des Institutsleiters>\\
<Straße und Hausnummer>\\
<Postleitzahl und Ort>\\[2ex]
Betreuer: <Name>\\[5ex]

% ----Tabelle der Praktikumsteilnehmer------------------------------------------
Auftragnehmer\\

\begin{tabular}{l<{\hspace{20mm}} l<{\hspace{30mm}}}\\
  Name                   &   E-Mail-Adresse\\      % Zeilenüberschift

  \hline                    % Linie unterhalb der Zeilenüberschrift

  %----Nachfolgend alle Namen und E-Mail-Adressen der Teilnehmer einfügen
  <Name des Teilnehmers> &  <E-Mail-Adresse>\\
  <Name des Teilnehmers> &  <E-Mail-Adresse>\\
  <Name des Teilnehmers> &  <E-Mail-Adresse>\\
  ...                    &  ...


\end{tabular}\\[2ex]

Braunschweig, <Abgabedatum>
\end{center}
\end{titlepage}
